\section{The paid menu response on the selected history}
\label{sec:finite-menu-history-handoff}
The centered-field argument requires the menu response and the ordinary
transition to belong to one history. We now identify their parameters and
prices. The analytic inputs are the funded path and Gaussian data used in
the preceding one-step estimates; this step makes no new choice of a
background transition or of a door transporting insertions.

Fix a finite history, adjacent scales $m+1,m$, and one point of its
interpolating path. Denote the fine representation by $R_f$, the result
by $R_c$, and the corresponding local coordinate change and insertion
transport by $\mathcal D$. The construction of the unmarked result fixes
all three. Use the same tuned representation, forced array and derivative
of the one-leaf response when constructing the menu terms of
Sections~\ref{sec:ex-lattice-rg-multisource} and
\ref{sec:ex-lattice-menu-jets}. Their quiet identity is therefore the
profile of $R_c$, and their singleton identity uses the hinge of precisely
$\mathcal D$.

Write $b_f,b_c\ge100$ for the support exponents, $N$ for the finite
occupancy bound, and $u\ge0$ for the menu shipment allowance. A menu offer
is measured at exponent $b_f-4$ and allowance $u$. Its marker weights are
positive, its modes nonempty, and the parent images of its named sites lie
in its crests. The shipment inequality of
Section~\ref{sec:ex-lattice-price-receipts}, together with the birth-region
conversion in Theorem~\ref{thm:ex-lattice-volant-supplier}, gives the
rooted majorants required for the paid component construction. This is
uniform over every finite menu satisfying the offer bound.

Let $a$ be the earned exponent and $e>0$ the additional reserve, so
$a\ge b_c+3$. Let $P_0,P_1\ge0$ be the ordinary and flagged component
budgets, $0\le\rho\le1/2$ the cluster ratio, and $E_0\ge0$ the marked
exponential cost. The field-price calculation gives
\begin{equation}\label{eq:menu-total-pledge}
 4e^{E_0}P_1+4\frac{P_0}{1-\rho}\le V_Q\le\tfrac12.
\end{equation}
Here the first inequality follows from the same component majorants used
in the cluster sums. To make its numerical content explicit, the total
pledge is $V_Q=10(P_0+E')$, with $P_1=E'e^{-E_0}$, $E'\ge0$.
Since $(1-\rho)^{-1}\le2$, its left side is at most
$4E'+8P_0\le10(E'+P_0)$. The upper bound $V_Q\le1/2$ is purchased
by increasing the entry coupling: both $P_0$ and $E'$ have the fixed
negative upper exponent of the component estimates. This is the same
arrival threshold used to pay their degree prices, so all inequalities
hold together.

The logarithmic response has the rooted bound
$2P_0/(1-\rho)$ at exponent $a+e$, by
\eqref{eq:ex-lattice-multi-lucent}. In particular it is at most $1$,
indeed at most $1/4$, by \eqref{eq:menu-total-pledge}.
Replacing its exponent by $b_c+1/2$ only decreases each nonnegative
summand, since $b_c+1/2\le a+e$ and all cover numbers are nonnegative.
The flagged response has the separate rooted bound
\[
                  t=2\exp\{-\mu\Delta_{G,L}(\beta)^2/16\}
\]
at exponent $b_c+1$. This is the second, larger-reserve response from
Section~\ref{sec:ex-lattice-rg-multisource}, paid by the two successive
absorptions of Section~\ref{sec:nonlinear-payment-selection}.
It is not deduced from the logarithmic bound.

The paid marker substitution and grouping now yield ordinary and flagged
families with exactly the ledgers
\eqref{eq:ex-lattice-volant-first-ledger} and
\eqref{eq:ex-lattice-volant-second-ledger}. Their supports, threshold
shells and quiet and singleton identities are those of
Theorem~\ref{thm:ex-lattice-volant-supplier}. The full push identity
retains the exponential of the ordinary polynomial and the complete
flagged gas. In particular no linearity of that full expression in the
marker families is used. The derivative construction fixes the tuning
law, and the retained regular and forced slots provide the input labels
used by the offer. Thus the response is available for the same door
$\mathcal D$ and the same pair $R_f,R_c$.

\subsection{Identifying the centered transfer allowances}
Let $\widehat\beta_m$ be the selected running coupling and $L$ the
block side at this step. Their exact relation is
\[
 \widehat\beta_{m+1}
       =\beta_F(\widehat\beta_m,L)
       =\widehat\beta_m+s_G\log L.
\]
The fine and coarse radii are correspondingly
$B_f=B(\widehat\beta_{m+1})$ and $B_c=B(\widehat\beta_m)$.
Take the offer allowance and the single-insertion price to be
\[
                   u=\widehat\beta_{m+1}^{-1},\qquad
                   \delta=\widehat\beta_{m+1}^{-2}.
\]
These are the tariff's station and probe prices; they are evaluated
at the same fine coupling. The tariff proved in
Section~\ref{sec:ex-lattice-price-receipts} says precisely that
$u\ge0$, $0<\delta\le1$, and
\begin{align*}
 T+e^TT&\le u,&T&=N\mathfrak s(B_f)\delta+4\delta,\\
 (2B_c+1)^4(t+N\mathfrak s(B_c)\delta)&\le\tfrac12,&
 e^{t+N\mathfrak s(B_c)\delta}-1&\le1.
\end{align*}
The exponent margins are $b_f-4$ for the input offer, $b_c+1/2$ for
the ordinary output and $b_c+1$ for its flagged output. These are
exactly the margins in the local served-transfer proof of
Section~\ref{sec:local-transfer}. The finite wing and its singleton
coefficients are transported by $\mathcal D$, whose soundness is the
exact one-leaf transition already retained above.

Finally, the point of the interpolating path represents the fine history
entry at $(N,m+1)$, and its result represents the coarse entry at $(N,m)$.
Substituting these two equalities in the formulas just proved changes no
map or coefficient. It identifies the door with the history's chosen
adjacent-scale door. Hence the paid menu response and the three displayed
payment inequalities are the ones required in
Section~\ref{ex-centered-construction}. The remaining construction of
the funded nonlinear path is independent of this handoff; conditional on
that funded path and its retained slot and representation data, there is
no additional finite-menu or tariff hypothesis here.
